\documentclass{cris}

%% ---------- Paper metadata ----------
\title{Transaction Cost Erosion in Momentum Strategy Implementation:\\ A Practitioner's Perspective}
\author{J.~Kim\textsuperscript{1} \and H.~Park\textsuperscript{2}}
\volume{1}
\papernumber{4}
\paperyear{2025}
\papermonth{December}
\shorttitle{Transaction Cost Erosion in Momentum}
\shortauthors{Kim \& Park}
\keywords{momentum, transaction costs, implementation shortfall, factor investing}
\jelcodes{G11, G12, G14}
\correspondence{submit.inquiry@strategycritique.org}

\begin{document}
\maketitle

%% ---------- Abstract ----------
\begin{abstract}
We quantify the gap between theoretical momentum returns and those achievable after realistic transaction costs. Using a dataset of institutional execution records from 2010--2024, we decompose the implementation shortfall of standard cross-sectional momentum strategies into market impact, timing delay, and opportunity cost components. Our findings suggest that approximately 40--60\% of gross momentum premia dissipate under institutional-scale execution, with market impact alone accounting for the majority of the erosion. These results have direct implications for portfolio construction and capacity estimation.
\end{abstract}

%% ---------- Introduction ----------
\section{Introduction}

The momentum anomaly, first documented by \citet{jegadeesh1993}, remains one of the most studied phenomena in empirical asset pricing. Academic estimates of momentum premia typically range from 6--12\% annually for long-short portfolios formed on 12-month past returns with a one-month skip.

However, practitioners implementing these strategies face a reality that academic studies frequently understate: the cost of turning over a momentum portfolio at the frequency required to capture the premium. As \citet{novy2012} and \citet{frazzini2014} have shown, transaction costs can substantially reduce or eliminate the profitability of many anomalies.

This note contributes to the growing literature on implementation costs by providing a decomposition of momentum execution costs using proprietary institutional data. Unlike prior studies that rely on estimated costs or TAQ-derived measures, our analysis draws on actual fill records from a systematic equity allocator.

%% ---------- Literature ----------
\section{Related Literature}

The tension between gross and net factor returns has received increasing attention. \citet{fama1993} established the standard factor framework, but their portfolio construction implicitly assumes frictionless rebalancing. \citet{daniel2016} documented the role of momentum crashes, which compound the cost problem by concentrating turnover in adverse market conditions.

Our work extends these findings by providing granular cost decomposition at the execution level --- a perspective available only to practitioners with access to implementation data.

%% ---------- Methodology ----------
\section{Methodology}

\subsection{Portfolio Construction}

We construct a standard cross-sectional momentum strategy following the specification of \citet{jegadeesh1993}: at the end of each month, stocks in the MSCI World Index are ranked by their cumulative return over months $t{-}12$ through $t{-}2$. The top decile forms the long portfolio; the bottom decile forms the short portfolio.

The portfolio is rebalanced monthly. Gross returns are computed assuming end-of-day closing prices. Net returns incorporate the actual execution costs described below.

\subsection{Cost Decomposition}

We decompose total implementation shortfall $S$ into three components:
\begin{equation}
  S = \underbrace{C_{\text{impact}}}_{\text{market impact}}
    + \underbrace{C_{\text{delay}}}_{\text{timing delay}}
    + \underbrace{C_{\text{opp}}}_{\text{opportunity cost}}
\end{equation}
where market impact $C_{\text{impact}}$ is measured as the difference between the volume-weighted average execution price and the arrival price; timing delay $C_{\text{delay}}$ captures the drift between decision time and order submission; and opportunity cost $C_{\text{opp}}$ reflects unfilled portions of target orders.

%% ---------- Results ----------
\section{Results}

Table~\ref{tab:costs} summarizes the annualized cost decomposition across our sample period.

\begin{table}[ht]
  \centering
  \caption{Annualized implementation shortfall decomposition for momentum strategy, 2010--2024. All figures in basis points.}
  \label{tab:costs}
  \begin{tabular}{lrrr}
    \toprule
    Component & Mean & Median & Std.~Dev. \\
    \midrule
    Market impact     & 187 & 172 & 64 \\
    Timing delay      &  43 &  38 & 21 \\
    Opportunity cost  &  52 &  41 & 35 \\
    \midrule
    Total shortfall   & 282 & 251 & 89 \\
    \bottomrule
  \end{tabular}
\end{table}

The mean total implementation shortfall of 282 basis points annually represents a substantial fraction of the gross momentum premium, which averaged approximately 640 basis points over the same period. Market impact alone accounts for 66\% of total costs, consistent with the high turnover and concentrated trading patterns inherent to momentum rebalancing.\footnote{Turnover rates for decile-based momentum strategies typically exceed 200\% annually on each leg of the portfolio.}

\subsection{Conditional Cost Patterns}

Implementation costs are not uniformly distributed across market regimes. During momentum crash episodes --- periods identified following the methodology of \citet{daniel2016} --- total shortfall increases by approximately 40\%, driven primarily by elevated market impact during forced liquidation of losing positions.

%% ---------- Discussion ----------
\section{Discussion}

Our findings carry several implications for practitioners:

\begin{enumerate}
  \item Gross momentum returns overstate implementable returns by approximately 44\% in our sample, and by more than 60\% during stress periods.
  \item Market impact, not commission or delay, is the dominant cost driver. This suggests that capacity management and execution algorithm design are first-order concerns for momentum allocators.
  \item The gap between academic estimates and implementable returns is not a fixed constant but varies with market conditions, making static cost adjustments inadequate for realistic performance assessment.
\end{enumerate}

These observations reinforce the value of practitioner-oriented research that complements theoretical analysis with implementation evidence.

%% ---------- Conclusion ----------
\section{Conclusion}

This note demonstrates that the momentum premium, while robust in gross terms, is substantially eroded by transaction costs at institutional scale. We provide a granular decomposition that practitioners can use to calibrate their own capacity estimates and execution benchmarks. Future work should extend this analysis to alternative momentum definitions, including residual and industry-adjusted momentum, which may exhibit different cost profiles.

%% ---------- References ----------
\section{References}
\bibliography{cris-sample}

\end{document}
